\subsection*{扩展数系，也扩展乘法能表达的动作}
实数乘法能沿数轴缩放或反向，却无法表示平面中的四分之一圈。引入 $\mathrm i$ 并规定 $\mathrm i^2=-1$，可以把实数对写成 $a+b\mathrm i$。乘以 $\mathrm i$ 得到 $-b+a\mathrm i$，正是把平面坐标转九十度。

两次四分之一圈合为半圈，所以 $\mathrm i^2=-1$ 的几何解释与代数规则一致。这个扩展不是让原有规则随意失效，而是在保留加法乘法基本规律的同时，让原本无实数解的方程得到新的表示。

\subsection*{同一个乘法有两种阅读方式}
直角坐标形式适合展开：按分配律相乘，再把 $\mathrm i^2$ 换成负一。极坐标形式适合看整体动作：模长相乘、角度相加。对 $(1+\mathrm i)^2$，展开得 $2\mathrm i$；几何上是两次四十五度旋转，模长 $\sqrt2$ 也相乘为二。

求倒数则需要撤销这两个动作：模长取倒数、角度取负。乘共轭能在坐标中实现同一件事，因为分母变成模长平方。请把 $1/(1+\mathrm i)$ 先画成反向旋转加缩小，再算出 $(1-\mathrm i)/2$。

复数有自然的长度和角度，却没有与实数完全相同的有序大小关系。谈“更大”时要说明是模长、实部还是别的量。下一节指数语言会把这套旋转与缩放规则进一步统一。
